\paperchapter{September 2013}{September 2013}{A}{60}{GSET 2013 (I)}
\keytable{\kcell{1}{A} & \kcell{2}{C} & \kcell{3}{C} & \kcell{4}{B} & \kcell{5}{C} & \kcell{6}{B} & \kcell{7}{D} & \kcell{8}{B} & \kcell{9}{B} & \kcell{10}{B} \\
\kcell{11}{B} & \kcell{12}{A} & \kcell{13}{D} & \kcell{14}{D} & \kcell{15}{C} & \kcell{16}{A} & \kcell{17}{C} & \kcell{18}{D} & \kcell{19}{D} & \kcell{20}{C} \\
\kcell{21}{C} & \kcell{22}{C} & \kcell{23}{A} & \kcell{24}{D} & \kcell{25}{C} & \kcell{26}{B} & \kcell{27}{C} & \kcell{28}{B} & \kcell{29}{A} & \kcell{30}{C} \\
\kcell{31}{A} & \kcell{32}{B} & \kcell{33}{B} & \kcell{34}{A} & \kcell{35}{B} & \kcell{36}{C} & \kcell{37}{C} & \kcell{38}{B} & \kcell{39}{C} & \kcell{40}{B} \\
\kcell{41}{C} & \kcell{42}{C} & \kcell{43}{C} & \kcell{44}{D} & \kcell{45}{B} & \kcell{46}{D} & \kcell{47}{C} & \kcell{48}{A} & \kcell{49}{A} & \kcell{50}{A} \\
\kcell{51}{C} & \kcell{52}{C} & \kcell{53}{D} & \kcell{54}{B} & \kcell{55}{A} & \kcell{56}{D} & \kcell{57}{A} & \kcell{58}{B} & \kcell{59}{A} & \kcell{60}{C}}
\begin{question}{1}
{Teaching means to help students\par}
\optsfour{how to think}{what to think}{why to think}{when to think}
\answer{A}{how to think}
\begin{explain}
Good teaching develops the learner's own thinking ability. It teaches students \textbf{how to think}: reasoning, analysis and problem-solving. Telling them \emph{what} to think is indoctrination.
\end{explain}
\end{question}

\begin{question}{2}
{Which of the following is not true when we say that `An educational institution is a laboratory that provides experiences'?\par}
\optsone{in absorbing and producing knowledge}{in getting life's adventure}{for practical economics and self-development}{for actually completing the course}
\answer{C}{for practical economics and self-development}
\begin{explain}
A school or college as a ``laboratory of experiences'' helps learners absorb and produce knowledge, experience life's adventure, and develop themselves. Per the official key, ``for \textbf{practical economics} and self-development'' is \emph{not} a correct description of that idea.
\end{explain}
\end{question}

\begin{question}{3}
{Regularity in class attendance can be ensured by\par}
\optsone{organising cultural programmes in the college}{imposing fine from students}{having frequent interaction with parents and students}{having strict discipline code in the college}
\answer{C}{having frequent interaction with parents and students}
\begin{explain}
Regular attendance is best ensured by \textbf{frequent interaction with parents and students}: understanding the reasons for absence, motivating students and involving parents.
\why{Cultural programmes, fines and strict codes are superficial or coercive.}
\end{explain}
\end{question}

\begin{question}{4}
{The purpose of using teaching aids in the classroom is to\par}
\optstwo{motivate students for learning}{concretise abstract concepts}{simplify teaching}{use imagination of students}
\answer{B}{concretise abstract concepts}
\begin{explain}
Teaching aids make \textbf{abstract concepts concrete} (``from abstract to concrete'') through models, charts, videos and demonstrations, so that they are easier to understand and remember.
\end{explain}
\end{question}

\begin{question}{5}
{A college teacher will really help students if she\par}
\optsone{explains at length}{is objective in her evaluation}{encourages students to learn on their own}{covers syllabus completely in the class}
\answer{C}{encourages students to learn on their own}
\begin{explain}
A college teacher helps most by \textbf{encouraging students to learn on their own}, which develops self-directed, lifelong learners.
\end{explain}
\end{question}

\begin{question}{6}
{Which of the following is a more desirable outcome of teaching in higher education?\par}
\optsone{Increase in student scores}{Enhanced independent thinking of students}{Higher pass percentage}{Increase in number of students who opt for the subject}
\answer{B}{Enhanced independent thinking of students}
\begin{explain}
The most desirable outcome of higher education is \textbf{enhanced independent thinking} of students, not scores or pass percentages.
\end{explain}
\end{question}

\begin{question}{7}
{Which of the following kinds of research studies does not require sampling?\par}
\optsfour{Ex-post Facto}{Experimental}{Survey}{Case study}
\answer{D}{Case study}
\begin{explain}
A \textbf{case study} is an in-depth study of a \emph{single} unit (a person, institution or event), so it needs no sampling from a population.
\end{explain}
\end{question}

\begin{question}{8}
{Which of the following pertains to scientific method ?\par}
\optstwo{experiential evidences}{hypothetico deduction}{a priori exploration}{belief in subjectivity of evidence}
\answer{B}{hypothetico deduction}
\begin{explain}
The scientific method is \textbf{hypothetico-deductive}: a hypothesis is formed, testable consequences are deduced from it, and these are checked against empirical evidence.
\end{explain}
\end{question}

\begin{question}{9}
{Opinion of 'facebook users' on reservation policy in education is to be studied. The most appropriate sampling technique for the purpose is:\par}
\optstwo{Stratified Random Sampling}{Snowball Sampling}{Cluster Sampling}{Quota Sampling}
\answer{B}{Snowball Sampling}
\begin{explain}
The key gives \textbf{snowball sampling}. Facebook users are a networked, hard-to-list population, and respondents can be reached through referral chains (friends of friends).
\end{explain}
\end{question}

\begin{question}{10}
{One's attitude towards knowledge should be:\par}
\optsone{to accept what is proved by empirical evidence}{to accept believing that knowledge leads to improvement}{to accept the final truth}{not to accept anything that is not personally verified}
\answer{B}{to accept believing that knowledge leads to improvement}
\begin{explain}
According to the key, one should accept knowledge \textbf{believing that knowledge leads to improvement}, i.e.\ with an open, progressive attitude rather than as final truth or with total scepticism.
\end{explain}
\end{question}

\begin{question}{11}
{Objectivity in research implies:\par}
\optstwo{Correct judgement of truth}{Findings consistent with reality}{Agreement in results}{Methodological sophistication}
\answer{B}{Findings consistent with reality}
\begin{explain}
\textbf{Objectivity} means the findings reflect reality and are free from the researcher's personal bias. The findings are \textbf{consistent with reality}.
\end{explain}
\end{question}

\begin{question}{12}
{The research that aims at immediate application is\par}
\optstwo{Action Research}{Empirical Research}{Conceptual Research}{Fundamental Research}
\answer{A}{Action Research}
\begin{explain}
\textbf{Action research} aims at solving immediate, practical problems in a specific setting and applying the results at once.
\end{explain}
\end{question}

\begin{question}{13}
{Find the odd one out\par}
\optsfour{CD ROM}{Hard Disk}{Pen Drive}{RAM}
\answer{D}{RAM}
\begin{explain}
CD-ROM, hard disk and pen drive are \emph{secondary (permanent) storage}. \textbf{RAM} is \emph{primary, volatile} memory, so it is the odd one out.
\end{explain}
\end{question}

\begin{question}{14}
{Radio and TV are more powerful and successful than the newspaper in communicating socially relevant messages because:\par}
\optsone{they are designed for such purposes}{they are capable of catching the attention of people}{they present a variety of programmes}{they can reach even the illiterate masses}
\answer{D}{they can reach even the illiterate masses}
\begin{explain}
Radio and TV use sound and pictures, so they \textbf{reach even the illiterate masses}, which newspapers cannot do.
\end{explain}
\end{question}

\begin{question}{15}
{One significant feature that enhances quality of teacher's communication in classroom is:\par}
\optsone{low noise level inside and outside classroom}{students' level of understanding}{teacher's ability to empathise with students}{accuracy in the content being communicated}
\answer{C}{teacher's ability to empathise with students}
\begin{explain}
A teacher who \textbf{empathises} with students understands their feelings and difficulties and adjusts the message accordingly. This raises the quality of classroom communication.
\end{explain}
\end{question}

\begin{question}{16}
{Full form of PDF is:\par}
\optstwo{Portable Document Format}{Portable Document File}{Permanent Document Format}{Permanent Document File}
\answer{A}{Portable Document Format}
\begin{explain}
\textbf{PDF} is \textbf{Portable Document Format} (Adobe, 1993).
\end{explain}
\end{question}

\begin{question}{17}
{The relation of data to information is that of\par}
\optstwo{raw material to its processing}{raw material to its quality}{raw material to a finished product}{both (A) and (B) above}
\answer{C}{raw material to a finished product}
\begin{explain}
Data are raw facts; information is processed, meaningful data. The relation is that of \textbf{raw material to a finished product}.
\end{explain}
\end{question}

\begin{question}{18}
{The set of procedures for handling the information transaction processing systems excludes\par}
\optsfour{calculation}{classification}{sorting}{conclusion}
\answer{D}{conclusion}
\begin{explain}
Transaction processing involves data operations: capturing, calculating, classifying, sorting, storing and summarising. Drawing a \textbf{conclusion} is a human decision-making activity, not a processing procedure.
\end{explain}
\end{question}

\begin{question}{19}
{To disrupt cellphone reception in a specific region jammers are used. This is a tool to obstruct communication at\par}
\optsfour{Receiver level}{Sender level}{Feedback level}{Media level}
\answer{D}{Media level}
\begin{explain}
Jammers block the radio signals in the transmission path, so they obstruct communication at the \textbf{media (channel) level}.
\end{explain}
\end{question}

\begin{question}{20}
{The potential of ICT is better utilized when used as a\par}
\optstwo{Teaching Supplement}{Teaching Support}{Learning Resource}{Learning Reinforcement}
\answer{C}{Learning Resource}
\begin{explain}
ICT is best used as a \textbf{learning resource} that learners access actively and independently, rather than merely as a teaching supplement.
\end{explain}
\end{question}

\begin{question}{21}
{Which TV channel is launched to provide technology based real classroom experiences through virtual mode to the distance education learners ?\par}
\optsfour{Gyan Darshan}{Eklavya}{Vyas}{Kisan}
\answer{C}{Vyas}
\begin{explain}
\textbf{Vyas} (the CEC--UGC higher-education channel, launched 2004) offers technology-based virtual classroom programmes for higher-education learners.
\why{Gyan Darshan (IGNOU, 2000) is a general educational channel; Eklavya is an engineering channel; Kisan is agricultural.}
\end{explain}
\end{question}

\begin{question}{22}
{A computer operating system that permits large number of students and teachers to run programs at the same time is known as\par}
\optstwo{Multiprogramming System}{Real Time System}{Multitasking System}{Time Sharing System}
\answer{C}{Multitasking System}
\begin{explain}
A \textbf{multitasking (multi-user)} operating system lets many users run programs at the same time.
\end{explain}
\end{question}

\begin{question}{23}
{What is the advantage of using database indexes over browsing the internet during research?\par}
\optsone{Information on database has been checked by the sources that publish it}{Database have wide variety of sources}{Access to relevant information is faster than the web browsing}{Databases are huge}
\answer{A}{Information on database has been checked by the sources that publish it}
\begin{explain}
Academic databases (indexes such as ERIC, Scopus or JSTOR) contain material that has been \textbf{screened and checked} (peer-reviewed or edited) by the publishers, unlike much of the open web. This is their chief advantage for research.
\end{explain}
\end{question}

\begin{question}{24}
{Which of the following acronyms would normally be used to describe unsolicited junk e-mails?\par}
\optsfour{WHAM}{DRAM}{CRAM}{SPAM}
\answer{D}{SPAM}
\begin{explain}
Unsolicited bulk junk e-mail is called \textbf{SPAM}.
\end{explain}
\end{question}

\begin{question}{25}
{Insert the missing number in the following sequence 240, \ldots{}, 120, 40, 10, 2\par}
\optsfour{120}{80}{240}{60}
\answer{C}{240}
\begin{explain}
Division by 1, 2, 3, 4, 5: $240\div1=\mathbf{240}$, $240\div2=120$, $120\div3=40$, $40\div4=10$, $10\div5=2$.
\end{explain}
\end{question}

\begin{question}{26}
{If BUT = 101 and NUT = 1301, then HUT = ?\par}
\optsfour{501}{701}{901}{1001}
\answer{B}{701}
\begin{explain}
In both BUT and NUT, ``UT'' gives the last two digits ``01''. The first letter gives the rest: B $=2$ gives $1$, N $=14$ gives $13$, i.e.\ position $-1$. H $=8$ gives $7$, so HUT $=\mathbf{701}$.
\end{explain}
\end{question}

\begin{question}{27}
{Complete the sequence B, D, G, I, L, N, \ldots{}\par}
\optsfour{O}{P}{Q}{R}
\answer{C}{Q}
\begin{explain}
The steps alternate $+2, +3$: B(2) $+2$ D $+3$ G $+2$ I $+3$ L $+2$ N $+3$ \textbf{Q}.
\end{explain}
\end{question}

\begin{question}{28}
{What will be the next letters in the following sequence ? CDBA, GHFE, KLJI, \ldots{}\par}
\optsfour{OPMN}{OPNM}{MNLK}{PQON}
\answer{B}{OPNM}
\begin{explain}
Each group is two consecutive letters followed by the two before them, in reverse: CD|BA, GH|FE, KL|JI, so the next is \textbf{OP|NM}.
\end{explain}
\end{question}

\begin{question}{29}
{What will be the next letters in the following sequence ? FEXW, GHVU, JITS, \ldots{}\par}
\optsfour{KLRQ}{LMRQ}{MLQR}{KLSR}
\answer{A}{KLRQ}
\begin{explain}
The first pair F E, G H, J I, \textbf{K L} alternates reversed/straight; the second pair X W, V U, T S, \textbf{R Q} runs backwards. The next term is \textbf{KLRQ}.
\end{explain}
\end{question}

\begin{question}{30}
{Q is shorter than P but taller than R. R is shorter than P but taller than A. W is shorter than A. If they stand in ascending order of their height, the order in which they are standing would be :\par}
\optsfour{ARWQP}{WAQRP}{WARQP}{WARPQ}
\answer{C}{WARQP}
\begin{explain}
P $>$ Q $>$ R $>$ A $>$ W. In ascending order: \textbf{W A R Q P}.
\end{explain}
\end{question}

\begin{context}
\textbf{Read the following text and answer the questions 31 to 36.}\par\smallskip
A nation is advanced in proportion of education and intelligence spread among the masses. The chief cause of India's ruin has been the monopolising of the whole education and intelligence of the land among handful of men. If we are to rise again, we shall have to do it by spreading education among the masses. The only service to be done for our lower classes is to give them education to develop their individuality. They are to be given ideas. Their eyes are to be opened to what is going on in the world around them, and then they will work out their own salvation. Every nation, every man and every woman must work out his own salvation. Give them ideas; that is the only help they require, and then the rest must follow as effect. Ours is to put the chemicals together; the crystallization comes in the law of nature. \hfill (Swami Vivekananda)
\end{context}
\begin{question}{31}
{Exposure to new ideas will empower masses as it enables them to understand\par}
\optsone{happenings in the world around}{that they need to open their eyes}{helps them to demand what is due to them}{that they can also learn}
\answer{A}{happenings in the world around}
\begin{explain}
``Their eyes are to be opened to \textbf{what is going on in the world around them}, and then they will work out their own salvation.''
\end{explain}
\end{question}

\begin{question}{32}
{Enhancing the proportion of education and intelligence among the masses can also be called\par}
\optstwo{National programme}{Mass education}{Development programme}{Educational programme}
\answer{B}{Mass education}
\begin{explain}
Spreading education and intelligence among the masses is \textbf{mass education}.
\end{explain}
\end{question}

\begin{question}{33}
{One major drawback faced by India post-independence was\par}
\optsone{unwillingness of masses to get educated}{concentration of education among smaller sections of population}{economic downfall of masses}{lack of ideas among the citizens}
\answer{B}{concentration of education among smaller sections of population}
\begin{explain}
``The chief cause of India's ruin has been the \textbf{monopolising of the whole education and intelligence} of the land among a handful of men.''
\end{explain}
\end{question}

\begin{question}{34}
{Upliftment of masses involves one of the following\par}
\optsone{everyone works out one's own salvation}{introduction of programmes for people's salvation by the Government}{freedom given to people to work independently}{everyone's eyes open while learning}
\answer{A}{everyone works out one's own salvation}
\begin{explain}
``\ldots and then they will work out their own salvation. \textbf{Every nation, every man and every woman must work out his own salvation.}''
\end{explain}
\end{question}

\begin{question}{35}
{Only education can help masses develop as\par}
\optsone{it makes them learn}{it enhances their individuality}{it is like putting several chemicals together}{it helps them crystallize their efforts}
\answer{B}{it enhances their individuality}
\begin{explain}
``The only service \ldots\ is to give them education \textbf{to develop their individuality}.''
\end{explain}
\end{question}

\begin{question}{36}
{The passage is in line with the present day view of education as a process of\par}
\optstwo{self help}{idea generation}{facilitation of learning}{crystallisation of ideas}
\answer{C}{facilitation of learning}
\begin{explain}
The teacher only ``puts the chemicals together'' (provides ideas) and the learner grows by himself or herself. This is the modern view of education as the \textbf{facilitation of learning}.
\end{explain}
\end{question}

\begin{question}{37}
{Which of the following is not a method of water conservation?\par}
\optstwo{Dual flush in toilets}{Groundwater recharging}{RO and UV treatment of water}{Rain water harvesting}
\answer{C}{RO and UV treatment of water}
\begin{explain}
Dual-flush toilets, groundwater recharge and rainwater harvesting \emph{conserve} water. \textbf{RO and UV treatment} \emph{purifies} water; RO in fact wastes a large share of the input water.
\end{explain}
\end{question}

\begin{question}{38}
{Ozone hole is considered harmful to us because\par}
\optsone{it can increase global warming}{it allows UV radiations to reach earth which may cause skin cancer}{it may allow entry of toxic gases from other planets}{all the gases will escape from this hole and we will be without oxygen}
\answer{B}{it allows UV radiations to reach earth which may cause skin cancer}
\begin{explain}
An ozone hole lets more \textbf{UV-B radiation} reach the Earth, causing skin cancer, cataracts and damage to plants and marine life.
\end{explain}
\end{question}

\begin{question}{39}
{Which of the following is not a major source of energy?\par}
\optsfour{fossil fuel}{biomass}{rains}{nuclear energy}
\answer{C}{rains}
\begin{explain}
Fossil fuels, biomass and nuclear power are major energy sources. \textbf{Rain} is not in itself a major energy source; hydropower uses stored water, but rainfall as such is not listed as a source.
\end{explain}
\end{question}

\begin{question}{40}
{Which of the following is an environmental problem associated with abandoned coal mines ?\par}
\optsone{thermal pollution of streams in the area}{acid drainage due to leaching of soil heaps by rain water}{air pollution caused by smog from ozone formation}{alkaline solutions that pollute streams}
\answer{B}{acid drainage due to leaching of soil heaps by rain water}
\begin{explain}
Abandoned coal mines produce \textbf{acid mine drainage}. Rainwater leaches pyrite (FeS$_2$) in spoil heaps, forming sulphuric acid and dissolving heavy metals that pollute streams.
\end{explain}
\end{question}

\begin{question}{41}
{Climate change may pose health related problems because of\par}
\optsone{malnutrition due to lack of food}{diseases due to increased pests}{diseases due to increased population of insect carriers}{failure of existing medical facility to control any epidemic}
\answer{C}{diseases due to increased population of insect carriers}
\begin{explain}
Warming extends the range and breeding season of \textbf{disease vectors} such as mosquitoes (malaria, dengue), increasing vector-borne diseases.
\end{explain}
\end{question}

\begin{question}{42}
{Which of the following is least likely to be an effect of global warming?\par}
\optsone{Extinction of some species}{Increase in sea water level}{Decreased hydrocarbon in plants}{Loss of fertile delta regions for agriculture}
\answer{C}{Decreased hydrocarbon in plants}
\begin{explain}
Species extinction, sea-level rise and loss of deltas are expected effects of global warming. ``\textbf{Decreased hydrocarbon in plants}'' is not a recognised effect.
\end{explain}
\end{question}

\begin{question}{43}
{To ensure quality in higher education, the NAAC provides\par}
\optstwo{Norms and Regulations}{Rules and Procedures}{Indicators of quality}{Case studies of good institutions}
\answer{C}{Indicators of quality}
\begin{explain}
NAAC evaluates institutions on criteria and key \textbf{indicators of quality} (7 criteria: curriculum, teaching--learning, research, infrastructure, student support, governance, values and best practices).
\end{explain}
\end{question}

\begin{question}{44}
{Which of the following is a regulatory body of technical education in India?\par}
\optsfour{NCTE}{MCI}{AIU}{AICTE}
\answer{D}{AICTE}
\begin{explain}
The \textbf{AICTE} regulates technical education.
\why{NCTE regulates teacher education; MCI regulated medical education; AIU is an association of universities, not a regulator.}
\end{explain}
\end{question}

\begin{question}{45}
{Distance learning system involves studying courses through\par}
\optsone{printed learning material in the college library}{printed as well as other media at one's place of living}{various methods of learning in college}{one's own efforts in the college without attending classes}
\answer{B}{printed as well as other media at one's place of living}
\begin{explain}
In distance learning the learner studies \textbf{printed as well as other media (audio, video, online) at his or her own place}, without regular attendance at an institution.
\end{explain}
\end{question}

\begin{question}{46}
{Universities or Higher Education Institutions like to get accredited by NAAC because\par}
\optsone{degrees get accepted in the academic world}{status of the institution is established}{get relative placement in academic world}{strengths of university get streamlined}
\answer{D}{strengths of university get streamlined}
\begin{explain}
Per the official key, accreditation helps the institution \textbf{streamline its strengths}: self-study and peer review identify strengths and weaknesses for continuous improvement.
\end{explain}
\end{question}

\begin{question}{47}
{Continuous and comprehensive evaluation is better than summative evaluation because it\par}
\optsone{keeps students busy all through the year}{gives a progressive picture of students' achievement}{helps in improving students' competencies in different areas}{would improve the annual result of the class}
\answer{C}{helps in improving students' competencies in different areas}
\begin{explain}
Continuous and comprehensive evaluation assesses learners regularly across scholastic and co-scholastic areas and feeds the results back into teaching. Its real advantage over one-time summative evaluation is that it \textbf{helps in improving students' competencies in different areas} through timely remediation.
\end{explain}
\end{question}

\begin{question}{48}
{Several experts criticise semester system in universities. Which of the following statements best endorses their view?\par}
\optsone{It does not provide sufficient time for deeper knowledge}{It increases the burden of students and teachers}{It controls free thinking of students.}{The focus is more on earning grades and marks than learning}
\answer{A}{It does not provide sufficient time for deeper knowledge}
\begin{explain}
The main criticism of the semester system is that each semester is short, so there is \textbf{not enough time for deep learning}; teaching becomes hurried and examination-driven.
\end{explain}
\end{question}

\begin{context}
The following data provide the information of the percentage of outgoing students from an academic institution after completing their study. Answer the question nos.~49 to 51 based on these data.
\begin{center}\small
\begin{tabular}{lccc}\toprule \textbf{Percentage of students} & \textbf{2010} & \textbf{2011} & \textbf{2012}\\\midrule
By campus recruitment & 33 & 37 & 39\\ For further study & 38 & 40 & 41\\ Unemployed & 4 & 3 & 2\\ To other activities & 25 & 20 & 18\\\midrule
Total no.\ of students & 1000 & 1100 & 1200\\\bottomrule
\end{tabular}
\end{center}
\end{context}
\begin{question}{49}
{What is the average percentage of students per year employed through campus interviews?\par}
\optsfour{$36\frac{17}{33}$}{37}{38}{35}
\answer{A}{$36\frac{17}{33}$}
\begin{explain}
Students placed: $330+407+468=1205$ out of $1000+1100+1200=3300$.
\[\frac{1205}{3300}\times100=36.515\ldots=36\tfrac{17}{33}\%\]
\end{explain}
\end{question}

\begin{question}{50}
{How many students choose other activities in 2012?\par}
\optsfour{216}{210}{250}{180}
\answer{A}{216}
\begin{explain}
$18\%$ of $1200=\mathbf{216}$.
\end{explain}
\end{question}

\begin{question}{51}
{What is the increase in the number of students going for further studies during 2010-2012?\par}
\optsfour{300}{30}{112}{110}
\answer{C}{112}
\begin{explain}
Further study: 2010 $=38\%\times1000=380$; 2012 $=41\%\times1200=492$. Increase $=\mathbf{112}$.
\end{explain}
\end{question}

\begin{context}
\textbf{Answer Q.~Nos.~52 to 54 based on the following chart.} The pie chart below shows various types of residences in a housing colony. In all there are 150 residential houses.
\begin{center}
\begin{tikzpicture}\pie[radius=1.4,text=pin,font=\scriptsize,color={brand!15,brand!35,brand!10,brand!45,brand!25}]{30/Flats,20/Bungalows,12/Independent houses,18/Row type,20/Semi-detached}\end{tikzpicture}
\end{center}
\end{context}
\begin{question}{52}
{The 20\% of the total bungalows are white in colour. How many are there?\par}
\optsfour{10}{8}{6}{12}
\answer{C}{6}
\begin{explain}
Bungalows $=20\%$ of $150=30$; white $=20\%$ of $30=\mathbf{6}$.
\end{explain}
\end{question}

\begin{question}{53}
{How many residences are flat type ?\par}
\optsfour{20}{30}{35}{45}
\answer{D}{45}
\begin{explain}
Flats $=30\%$ of $150=\mathbf{45}$.
\end{explain}
\end{question}

\begin{question}{54}
{From among the 12\% independent houses, 50\% are facing east direction. How many are they?\par}
\optsfour{12}{9}{6}{10}
\answer{B}{9}
\begin{explain}
Independent houses $=12\%$ of $150=18$; facing east $=50\%$ of $18=\mathbf{9}$.
\end{explain}
\end{question}

\begin{question}{55}
{In the following if the smallest circle represents candidates clearing GSET exam, bigger circle represents candidates clearing only Paper-I of the GSET exam and the biggest circle represents the candidates who failed in the GSET exam, then choose the diagram which correctly represents this relationships.\par}
\optsfour{\tikz[baseline=-0.5ex,scale=0.3]{\draw (0,0) circle(1.1); \draw (0,0) circle(0.75); \draw (2.2,0) circle(0.6);}}{\tikz[baseline=-0.5ex,scale=0.3]{\draw (0,0) circle(1); \draw (2.2,0) circle(0.8); \draw (3.8,0) circle(0.45);}}{\tikz[baseline=-0.5ex,scale=0.3]{\draw (0,0) circle(1.1); \draw (2.4,0) circle(0.8); \draw (2.4,0) circle(0.4);}}{\tikz[baseline=-0.5ex,scale=0.3]{\draw (0,0) circle(1.2); \draw (0,0) circle(0.85); \draw (0,0) circle(0.5);}}
\answer{A}{\tikz[baseline=-0.5ex,scale=0.3]{\draw (0,0) circle(1.1); \draw (0,0) circle(0.75); \draw (2.2,0) circle(0.6);}}
\begin{explain}
Per the key, (A) is correct. Candidates who clear GSET have necessarily cleared Paper I, so the \textbf{smallest circle lies inside the bigger circle}. Candidates who \textbf{failed} form a separate group, the \textbf{biggest circle, disjoint} from both.
\end{explain}
\end{question}

\begin{question}{56}
{Consider the following statement followed by two assumptions (i) and (ii). Choose the correct answer.\\
\textbf{Statement:} ``You are hereby appointed as a programmer with a probation period of one year and your performance will be reviewed at the end of the period for confirmation.'' --- A line in an appointment letter.\\
\textbf{Assumptions:}\\
(i) The performance of an individual is generally not known at the time of appointment offer.\\
(ii) Generally an individual tries to prove his/her worth in the probation period.\par}
\optsone{Only Assumption (i) is implicit.}{Only Assumption (ii) is implicit.}{None of the assumptions (i) and (ii) are implicit.}{Both the Assumptions (i) and (ii) are implicit.}
\answer{D}{Both the Assumptions (i) and (ii) are implicit.}
\begin{explain}
A probation period is used \emph{because} an appointee's real performance is not known at the time of appointment (i), and because people generally try to prove their worth during probation (ii). \textbf{Both} assumptions are implicit.
\end{explain}
\end{question}

\begin{question}{57}
{Of the following statements, there are two statements both of which cannot be false but both can be true. Which are these two statements?\\
(i) All the students wish to complete their study.\\
(ii) Some students do not like Mathematics.\\
(iii) No student wishes to complete the study.\\
(iv) No student likes Mathematics.\par}
\optsfour{(i) and (ii)}{(iii) and (iv)}{(i) and (iii)}{(ii) and (iv)}
\answer{A}{(i) and (ii)}
\begin{explain}
The official key gives \textbf{(i) and (ii)}. The pair ``All students wish to complete their study'' and ``Some students do not like Mathematics'' can obviously both be true together, and the setters treat them as the pair that cannot both be false.
\tip{In formal logic, two propositions that cannot both be false but can both be true are \emph{sub-contraries} (``Some S are P'' and ``Some S are not P'').}
\end{explain}
\end{question}

\begin{question}{58}
{Evaluate the following statements in terms of whether each is a Fact, Opinion, Prejudice or Advice by choosing correct alternative.\\
(i) Women are not suitable for police service.\\
(ii) In democracy no party should be in power for too long.\\
(iii) Have proper rest even during examination time.\\
(iv) Obesity is a risk factor for coronary heart diseases.\\
Options list the statements in the order: \textbf{Opinion -- Prejudice -- Fact -- Advice}.\par}
\optstwo{(ii) (iii) (iv) (i)}{(ii) (i) (iv) (iii)}{(iii) (i) (iv) (ii)}{(ii) (iv) (iii) (i)}
\answer{B}{(ii) (i) (iv) (iii)}
\begin{explain}
(ii) is an \textbf{opinion}; (i) is a \textbf{prejudice} (a biased generalisation about women); (iv) is a scientific \textbf{fact}; (iii) is \textbf{advice}. The order is (ii), (i), (iv), (iii), which is option (B).
\end{explain}
\end{question}

\begin{question}{59}
{The question below is followed by two arguments numbered (i) and (ii). Decide which argument is strong.\\
\textbf{Question:} Should India encourage exports when most things are insufficient for internal use itself?\\
\textbf{Arguments:}\\
(i) Yes. We have to earn foreign exchange to pay at least for our defence and other necessary imports.\\
(ii) No. Even selective encouragement will lead to shortages.\par}
\optsone{Only Argument (i) is strong.}{Only Argument (ii) is strong.}{Both the Arguments (i) and (ii) are strong.}{None of the arguments (i) and (ii) are strong.}
\answer{A}{Only Argument (i) is strong.}
\begin{explain}
(i) is \textbf{strong}: foreign exchange is essential for necessary imports such as defence equipment and oil. (ii) is weak: it assumes, without basis, that \emph{any} selective export will cause shortages.
\end{explain}
\end{question}

\begin{question}{60}
{Consider the following two statements and choose the correct answer.\\
\textbf{Statement (i):} There has been a high increase in the incidents of atrocities against women in the country during past few years.\\
\textbf{Statement (ii):} The police authority has been unable to nab the culprits committing crime against women.\par}
\optsone{Statement (i) is cause and Statement (ii) is effect.}{Statement (ii) is cause and Statement (i) is effect.}{Statements (i) and (ii) are both independent of causes.}{Statements (i) and (ii) are both effects of some independent causes.}
\answer{C}{Statements (i) and (ii) are both independent of causes.}
\begin{explain}
The official key treats the two statements as \textbf{independent} of each other: the rise in crime has many social causes, and police failure is a separate issue. Neither is shown to cause the other.
\end{explain}
\end{question}
